% Formas lineales: el gradiente no es un vector (partición por clase: apunte/scripts/reorganizar_por_clase.py)
\chapter{Formas lineales: el gradiente no es un vector}\label{cap:formas-lineales}

\section{La idea}

\section{Construcciones, cuentas y ejemplos}

\subsection{Producto vectorial: unir dos puntos, cortar dos rectas}\label{sec:vectorial}

\textbf{Desarrollo.} Queremos una operación que, dados dos puntos \(p,q\), devuelva la recta
\(\ell=p\times q\) que los une. ¿Qué tiene que cumplir? Que \(p\) esté en esa recta y \(q\)
también:
\[
  (p\times q)\cdot p=0,\qquad (p\times q)\cdot q=0 .
\]
Por lo tanto la operación debe producir un vector \emph{ortogonal a sus dos argumentos}. Y
más: si ponemos dos veces el mismo punto no hay recta que determinar, así que \(p\times p=0\).
Estas exigencias, junto con la linealidad en cada argumento, \emph{obligan} a que la
operación sea, a menos de una constante, el producto vectorial de \(\R^3\) (en homogéneas
la constante no importa):
\[
  p\times q=\begin{pmatrix}p_2q_3-p_3q_2\\ p_3q_1-p_1q_3\\ p_1q_2-p_2q_1\end{pmatrix}.
\]

\begin{resultado}[Las dos propiedades que lo definen]
\begin{itemize}
  \item \textbf{Linealidad} en cada argumento: \((2p)\times q=2\,(p\times q)\),
  \((p+p')\times q=p\times q+p'\times q\).
  \item \textbf{Antisimetría}: \(p\times q=-\,q\times p\). Es equivalente a \(p\times p=0\):
  el único vector igual a su opuesto es el nulo.
\end{itemize}
Son exactamente las propiedades del \emph{área orientada} de un paralelogramo (duplicar un
lado duplica el área; cambiar el orden de los lados cambia la orientación; dos lados
paralelos no encierran área). Por eso \(\lVert p\times q\rVert\) es el área del paralelogramo
de lados \(p,q\).
\end{resultado}

Por simetría de la ecuación de incidencia, la misma operación resuelve el problema dual:
\emph{el punto donde se cortan dos rectas es} \(\ell\times m\). Dos rectas paralelas
\((a,b,c)\) y \((a,b,c')\) dan \(\ell\times m=(c'-c)\,(b,-a,0)^T\): un punto ideal, su
dirección común.

\subsection{Formas lineales: la recta como máquina de medir}\label{sec:formas-lineales}

{\small\vspace{1mm}\linkUnidad{06_formas_lineales}{Simulación interactiva: Unidad 8 — Formas lineales y dualidad}\par}
{\small\vspace{1mm}\linkAnimacion{01_dualidad_gradiente}{Animación: dualidad punto–gradiente}\par}


\textbf{Desarrollo} Hasta acá los
vectores fueron flechas: cosas que se suman y se estiran. Hay otro objeto, igual de simple,
que en los libros aparece tarde y que conviene presentar apenas definido el
espacio: la \emph{forma lineal}, una función que toma un vector y devuelve un número,
respetando sumas y escalas.

\paragraph{Primero en \(\R\).} Una forma lineal en la recta es multiplicar por un número:
\(\varphi(x)=2x\). Es una regla graduada: los puntos donde \(\varphi\) vale
\(0,1,2,3,\ldots\) están equiespaciados sobre la recta, a distancia \(\tfrac12\). No hay
nada más que saber. Al pasar al plano la regla sigue siendo la misma, pero ahora hay
direcciones que la regla no mide.

\paragraph{Un ejemplo con números chicos.} Sea
\[
  \varphi(x,y)=2x+3y .
\]
Con \(v=(1,0)\) da \(2\); con \((0,1)\) da \(3\); con \((1,1)\) da \(5\); con \((3,-2)\)
da \(0\). Es lineal: \(\varphi(v+w)=\varphi(v)+\varphi(w)\) y
\(\varphi(\lambda v)=\lambda\,\varphi(v)\). Y basta conocer \(\varphi\) en los dos vectores
de la base para conocerla en todos: \(\varphi(x,y)=x\,\varphi(1,0)+y\,\varphi(0,1)\).

\paragraph{Hecho geométrico: una regla graduada tirada en el plano.} Los puntos donde
\(\varphi\) vale lo mismo forman una recta, \(2x+3y=c\), y al variar \(c\) se obtiene una
familia de rectas paralelas equiespaciadas. La forma lineal \emph{es} esa familia: para
medir un punto, se mira qué recta de la familia pasa por él. La recta \(c=0\) pasa por el
origen y contiene la dirección \((3,-2)\), que es lo que la forma ``no ve'' (su núcleo).
El vector \((2,3)\) es perpendicular a todas las rectas de nivel y apunta hacia donde
\(\varphi\) crece.

\begin{center}
\begin{tikzpicture}[scale=.78]
  \begin{scope}
    \clip (-3.2,-3.2) rectangle (4.2,3.2);
    \foreach \c in {-6,-4,-2,0,2,4,6,8}
      \draw[teal,line width=.7pt] (-3.2,{(\c+6.4)/3})--(4.2,{(\c-8.4)/3});
    \draw[teal,line width=1.2pt] (-3.2,{6.4/3})--(4.2,{-8.4/3});
    \draw[axisline] (-3.2,0)--(4.2,0);
    \draw[axisline] (0,-3.2)--(0,3.2);
    \draw[rust,line width=1pt,-{Latex[length=2.4mm]}] (0,0)--(1,1.5)
      node[right,font=\small,labelback] {\((2,3)\)};
    \draw[tealdark,line width=1pt,-{Latex[length=2.4mm]}] (0,0)--(1.5,-1)
      node[below right,font=\small,labelback] {\((3,-2)\): \(\varphi=0\)};
    \fill[ink] (1,1) circle (2pt) node[above left,font=\small,labelback] {\((1,1)\): \(\varphi=5\)};
  \end{scope}
  \foreach \c/\y in {8/-0.13,6/-0.8,4/-1.47,2/-2.13}
    \node[right,font=\scriptsize,text=teal] at (4.2,\y) {\(c=\c\)};
  \foreach \c/\y in {0/2.13,-2/1.47,-4/0.8,-6/0.13}
    \node[left,font=\scriptsize,text=teal] at (-3.2,\y) {\(c=\c\)};
\end{tikzpicture}
\end{center}

\paragraph{Las selectoras y el espacio de las formas.} Las formas más básicas son las que
seleccionan una coordenada: \(\mathsf{x}(x,y)=x\) y \(\mathsf{y}(x,y)=y\). Cualquier otra
es combinación de ellas: \(\varphi=2\,\mathsf{x}+3\,\mathsf{y}\). Las formas lineales se
suman y se multiplican por escalares, y el resultado sigue siendo una forma lineal: forman
un espacio vectorial, el \emph{dual} del plano. En la base canónica una forma se describe
con dos números, igual que un vector, pero conviene escribirla como \emph{fila}:
\[
  \varphi(v)=\begin{pmatrix}2&3\end{pmatrix}\begin{pmatrix}x\\y\end{pmatrix}
  =(2,3)\cdot v .
\]
La fila por la columna es el producto interno con el vector \((2,3)\): la forma se
``disfraza'' de vector gracias al producto interno. En coordenadas homogéneas esto ya se
usó: una recta \(\ell=(a,b,c)\) es una forma lineal, un punto \(p\) es un vector, y
\(\ell\cdot p=0\) es la incidencia de la sección~\ref{sec:interno}.

\paragraph{La matriz como pila de formas.} Una transformación lineal del plano al plano
devuelve dos números por cada vector; cada uno de ellos es una forma lineal. Con
\[
  A=\begin{pmatrix}2&3\\1&-1\end{pmatrix},\qquad
  A\begin{pmatrix}1\\2\end{pmatrix}=\begin{pmatrix}2\cdot1+3\cdot2\\1\cdot1-1\cdot2\end{pmatrix}
  =\begin{pmatrix}8\\-1\end{pmatrix},
\]
la primera fila mide con la regla \(2x+3y\) y la segunda con la regla \(x-y\).
Geométricamente, aplicar \(A\) es leer dos reglas graduadas a la vez. El núcleo de \(A\)
es el conjunto de puntos donde las dos reglas marcan cero: la intersección de las rectas
\(2x+3y=0\) y \(x-y=0\), que es sólo el origen porque las rectas se cortan
(\(\det A=-5\neq0\)). Si en cambio \(B=\begin{pmatrix}2&3\\4&6\end{pmatrix}\), las dos
filas son la misma regla escalada, sus rectas de nivel cero coinciden, el núcleo es toda
la recta \(2x+3y=0\) y \(\det B=0\): es la sección~\ref{sec:nucleo} vista desde las
filas. Y el método del producto cruz para autovectores (sección~\ref{sec:autovalores})
busca justamente el vector que todas las reglas de \(A-\lambda I\) marcan en cero.

\paragraph{Cómo se transforma una forma cuando se transforma el vector.} Sea
\(A=\begin{pmatrix}1&1\\0&2\end{pmatrix}\) y \(v=(1,1)\). Entonces \(Av=(2,2)\) y
\(\varphi(Av)=2\cdot2+3\cdot2=10\). ¿Qué regla, aplicada al \(v\) original, da ese mismo
\(10\)? Se calcula la fila por la matriz:
\[
  (2,3)\,A=\bigl(2\cdot1+3\cdot0,\ 2\cdot1+3\cdot2\bigr)=(2,8),
  \qquad (2,8)\cdot(1,1)=10\ \checkmark
\]
El vector se multiplica por \(A\) a la izquierda (columna); la forma se multiplica por
\(A\) a la derecha (fila). Se resume como \(\varphi(A\vec v)=(\varphi A)(\vec v)\):
``una forma lineal se transforma del lado contrario que un vector''.

Si lo que se quiere es \emph{arrastrar la regla junto con la figura}, es decir, encontrar
la forma \(\psi\) cuyas rectas de nivel son las imágenes por \(A\) de las rectas de nivel
de \(\varphi\), hay que deshacer \(A\): \(\psi=\varphi A^{-1}\). Con
\(A^{-1}=\begin{pmatrix}1&-\tfrac12\\0&\tfrac12\end{pmatrix}\),
\[
  \psi=(2,3)\,A^{-1}=\bigl(2,\ -1+\tfrac32\bigr)=\bigl(2,\tfrac12\bigr).
\]
Verificación: el punto \((3,-2)\) está en la recta \(\varphi=0\); su imagen es
\(A(3,-2)=(1,-4)\), y \(\psi(1,-4)=2-2=0\): sigue en la recta de nivel cero de la regla
arrastrada. El punto \((3,0)\), con \(\varphi=6\), va a \((3,0)\) y \(\psi(3,0)=6\).
Hecho geométrico: al deformar el plano las rectas de nivel se deforman con él, siguen
siendo rectas paralelas equiespaciadas (es una afinidad), pero la separación y la
inclinación cambian, y por eso los coeficientes de la forma cambian ``al revés'' que los
de los vectores.

\paragraph{El gradiente es una forma lineal.} Sea \(f(x,y)=x^2+y\) en el punto \((1,1)\),
donde \(f=2\). Para un desplazamiento chico \((h,k)\),
\[
  f(1+h,1+k)=2+2h+k+h^2\approx 2+\underbrace{(2h+k)}_{\text{forma lineal}} .
\]
La derivada en \((1,1)\) es la regla \(2h+k\): la mejor aproximación lineal de \(f\) cerca
del punto. Su recta de nivel cero, \(2(x-1)+(y-1)=0\), es la tangente a la curva de nivel
\(x^2+y=2\) en \((1,1)\). Que ``el gradiente es perpendicular a la curva de nivel'' es la
misma figura de arriba: el vector \((2,1)\) es perpendicular a la familia de rectas de la
forma. Ir a lo chico (la forma lineal en un punto) para volver a lo grande (la superficie
completa) es lo que hace el capítulo~\ref{cap:derivada} con el plano tangente y la métrica.

\begin{idea}[Aplicación en Inteligencia Artificial: Support Vector Machines (SVM) como formas lineales]
Desarrollo complementario, desde la especialización en IA (CEIA); no es de la cátedra. Un
\emph{clasificador lineal} separa dos lotes de puntos del plano (dos clases, \(y_i=\pm1\))
con una regla \(\omega(x)=w_1x+w_2y+b\): la clase es el signo de \(\omega\). Es exactamente
una forma lineal, una regla graduada tirada sobre la nube: los puntos viven en \(V\), el
clasificador vive en \(V^{*}\), y en coordenadas homogéneas \((x,y,1)\) la traslación \(b\)
entra en la misma fila, \(\omega=(w_1,w_2,b)\), igual que una recta en la
sección~\ref{sec:interno}. La máquina de vectores soporte elige, entre todas las reglas
que separan, la que deja más \emph{margen}: pide \(\omega(x_i)\ge+1\) para una clase y
\(\le-1\) para la otra, y entre las que cumplen toma la de \(\|w\|\) mínima. Las tres
rectas de nivel \(-1\), \(0\), \(+1\) son el borde de un lado, la frontera y el borde del
otro; los puntos que tocan los bordes son los \emph{vectores soporte}. El ancho de la
franja entre las rectas \(-1\) y \(+1\) es \(2/\|w\|\), porque la separación entre
curvas de nivel consecutivas de una forma es \(1/\|w\|\) (cuanto más
largo el vector de la forma, más apretadas las curvas), y esa cuenta necesita el
producto interno: sin métrica hay regla, pero no hay ancho. Ejemplo con enteros: clase
\(+1\) en \((1,1)\) y \((2,0)\), clase \(-1\) en \((-1,-1)\) y \((-2,0)\). La frontera
de margen máximo es \(x+y=0\) con \(w=(\tfrac12,\tfrac12)\), \(b=0\): los cuatro puntos
dan \(\omega=\pm1\) exacto (los cuatro son vectores soporte) y el ancho es
\(2/\|w\|=2\sqrt2\); la recta vertical \(x=0\) también separa, pero con ancho \(2\). Lo
verificamos resolviendo el problema de minimización con \texttt{scipy.optimize}. La unidad
\texttt{06\_formas\_lineales} (apéndice~\ref{ap:notebooks}) lo dibuja con el helper
\texttt{svm\_margen\_duro\_2d} de \file{cb104.py}: mover un punto y ver cómo gira la regla.
\end{idea}

\begin{resultado}[Dualidad en pocas palabras]
\begin{itemize}
  \item Un vector es una columna; una forma lineal es una fila. Fila por columna es un
  número: la medición.
  \item Hecho geométrico: vector = flecha; forma = familia de rectas paralelas
  equiespaciadas (una regla). En \(\R^3\), planos paralelos.
  \item Una matriz es una pila de formas (sus filas) o una hilera de vectores (sus
  columnas). Ambas lecturas son útiles: las filas dan el núcleo, las columnas dan la imagen.
  \item Al transformar el plano con \(A\), los vectores van con \(A\) y las formas con
  \(A^{-1}\) (o con \(A\) del otro lado, si se pregunta ``qué regla mide lo mismo'').
  \item En homogéneas, los puntos son vectores y las rectas son formas; cambiarlos de rol
  es la dualidad de la geometría proyectiva.
\end{itemize}
\end{resultado}

\section{En profundidad}
